\subsection{自同态的迹}

设 C 为交换环，E 为 C-模。映射 \((x^{*}, x) \mapsto \langle x, x^{*} \rangle\)（从 \(E^{*} \times E\) 到 C）是 C-双线性的，因为对所有 \(\gamma \in C\)，\(\langle \gamma x, x^{*} \rangle = \gamma \langle x, x^{*} \rangle\) 且 \(\langle x, x^{*} \gamma \rangle = \langle x, x^{*} \rangle \gamma\)；我们由此得到一个典范 C-线性映射

\[
\tau : \mathbf {E} ^ {*} \otimes_ {\mathbf {c}} \mathbf {E} \rightarrow \mathbf {C}\tag{16}
\]

使得 \(\tau(x^{*} \otimes x) = \langle x, x^{*} \rangle\)（§3，第5号）。现在再设E是有限生成的投射C-模；于是第2号的典范同构 (11) 是C-模同构，因此我们可以通过转移结构定义一个典范线性形式 \(\mathrm{Tr} = \tau \circ \theta_{\mathbf{E}}^{-1}\)（在C-模 \(\operatorname{End}_{\mathbf{C}}(\mathbf{E})\) 上）。对所有 \(u \in \operatorname{End}_{\mathbf{C}}(\mathbf{E})\)，标量 \(\operatorname{Tr}(u)\) 称为自同态 \(u\) 的迹；每个 \(u \in \operatorname{End}_{\mathbf{C}}(\mathbf{E})\) 都可以（一般以无穷多种方式）写成 \(x \mapsto \sum_{i} \langle x, x_{i}^{*} \rangle y_{i}\) 的形式，其中 \(x_{i}^{*} \in \mathbf{E}^{*}\)、\(y_{i} \in \mathbf{E}_{i}\)（依据第2号命题2的推论）；则

\[
\operatorname{Tr} (u) = \sum_ {i} \left\langle y _ {i}, x _ {i} ^ {*} \right\rangle \quad (\text { cf. } \S 1 0, \text { no. } 1 1).\tag{17}
\]

由定义，

\begin{enumerate}
\def\labelenumi{(\arabic{enumi})}
\setcounter{enumi}{17}
\tightlist
\item
\end{enumerate}

\[
\operatorname{Tr} (u + v) = \operatorname{Tr} (u) + \operatorname{Tr} (v)\tag{19}
\]

\[
\operatorname{Tr} (\gamma u) = \gamma \operatorname{Tr} (u)
\]

对 \(u, v\)（\(\operatorname{End}_{\mathbf{C}}(\mathbf{E})\) 中）及 \(\gamma \in \mathbf{C}\) 成立。此外：

命题3. 设 C 为交换环，E、F 为两个有限生成投射 C-模，且 \(u: \mathbf{E} \to \mathbf{F}\) 和 \(v: \mathbf{F} \to \mathbf{E}\) 为两个线性映射；则

\[
\operatorname{Tr} (v \circ u) = \operatorname{Tr} (u \circ v).\tag{20}
\]

映射 \((u, v) \mapsto \operatorname{Tr}(u \circ v)\) 与 \((u, v) \mapsto \operatorname{Tr}(v \circ u)\)（从

\[
\operatorname{Hom} _ {\mathbf {c}} (\mathbf {E}, \mathbf {F}) \times \operatorname{Hom} _ {\mathbf {c}} (\mathbf {F}, \mathbf {E})
\]

到C）都是 C-双线性的；因此只需在 u 形如 \(x \mapsto \langle x, a^{*} \rangle b\) 且 v 形如 \(y \mapsto \langle y, b^{*} \rangle a\)（其中 \(a \in E\)、\(a^{*} \in E^{*}\)、\(b \in F\)、\(b^{*} \in F^{*}\)）的情形验证 (20)。但此时 \(v \circ u\) 是映射 \(x \mapsto \langle x, a^{*} \rangle \langle b, b^{*} \rangle a\)，而 \(u \circ v\) 是映射 \(y \mapsto \langle y, b^{*} \rangle \langle a, a^{*} \rangle b\)。公式 (17) 表明 (20) 两端的值都等于 \(\langle a, a^{*} \rangle \langle b, b^{*} \rangle\)。

推论. 若 \(u_1, \ldots, u_p\) 是 \(\mathbf{E}\) 的自同态，则

\[
\operatorname{Tr} \left(u _ {1} \circ u _ {2} \circ \dots \circ u _ {p}\right) = \operatorname{Tr} \left(u _ {i} \circ u _ {i + 1} \circ \dots \circ u _ {p} \circ u _ {1} \circ \dots \circ u _ {i - 1}\right)
\]

对所有 \(1 \leqslant i \leqslant p\) 成立（“迹在循环置换下的不变性”）。

只需将 (20) 应用于乘积

\[
\left(u _ {1} \circ u _ {2} \circ \dots \circ u _ {i - 1}\right) \circ \left(u _ {i} \circ u _ {i + 1} \circ \dots \circ u _ {p}\right).
\]

注意另一方面，对于 E 的三个自同态 u、v、w，

\[
\operatorname{Tr} (u \circ v \circ w) = \operatorname{Tr} (u \circ w \circ v)
\]

不一定成立。

